\documentclass[11pt,a4paper]{article}
\usepackage[utf8]{inputenc}
\usepackage[T1]{fontenc}
\usepackage[french]{babel}
\usepackage{geometry}
\usepackage{hyperref}
\usepackage{enumitem}
\geometry{margin=2.3cm}

\title{Global Regularity and Spectral Confinement\\ via the Dissipative Invariant \\ \large Version augmentee fractale et chiastique}
\author{Sovereign Research Collective}
\date{March 2026}

\begin{document}
\maketitle

\section*{A. Ouverture}
Ce document augmente le manuscrit source en le replissant selon l'axe chiastique A-B-C-X-C'-B'-A'.

\section*{B. Expansion}
Les noeuds structurels extraits du manuscrit sont :
\begin{itemize}[leftmargin=1.2cm]
\item The Axiomatic of Invariant Margins
\item Global Regularity for 3D Navier--Stokes
\item The Spectral Problem and Riemann's Conjecture
\item Mass Gap and Yang--Mills Stability
\item The Birch and Swinnerton-Dyer (BSD) Conjecture
\item The Hodge Conjecture: Cohomological Regularity
\end{itemize}

\section*{C. Matiere}
Extrait interpretable du manuscrit d'origine :
\begin{quote}
documentclass[11pt,a4paper]{amsart} % ----------------------- PACKAGES ------------------------------------ usepackage[utf8]{inputenc} usepackage[T1]{fontenc} usepackage{amsmath, amssymb, amsthm} usepackage{geometry} usepackage{cite} usepackage{hyperref} usepackage{mathrsfs} usepackage{enumitem} geometry{margin=1in} hypersetup{ colorlinks=true, linkcolor=blue, citecolor=red, urlcolor=blue } % ----------------------- THEOREMS ------------------------------------- newtheorem{theorem}{Theorem}[section] newtheorem{lemma}[theorem]{Lemma} newtheorem{definition}[theorem]{Definition} newtheorem{remark}[theorem]{Remark} newtheorem{proposition}[theorem]{Proposition} newtheorem{corollary}[theorem]{Corollary} newtheorem{axiom}[theorem]{Axiome} % ----------------------- MACROS -------------------------------------- newcommand{alphacoeff}{ensuremath{alpha}} newcommand{betacoeff}{ensuremath{beta}} newc
\end{quote}

\section*{X. Centre}
Le centre chiastique est la stabilite dissipative comme principe de fermeture du manuscrit.

\section*{C'. Retour}
Le retour referme le texte sur sa fonction : preuve, seuil, verification ou acte d'unification.

\section*{B'. Miroir}
Le miroir fait correspondre architecture du manuscrit, autorite de son regime et lisibilite de sa trajectoire.

\section*{A'. Cloture}
La cloture ne remplace pas la source ; elle l'installe comme figure fractale augmentee dans l'arche souveraine.

\section*{Metadonnees}
Source : \verb|_RELEASES/GOLDEN_MASTER_PHASE_III/Sovereign_Millennium_Dissipative_Stability_Proof.tex|\\
Titre detecte : \texttt{Global Regularity and Spectral Confinement\\ via the Dissipative Invariant}\\
Auteur detecte : \texttt{Sovereign Research Collective}\\
Date detectee : \texttt{March 2026}

\end{document}
